\begin{afterword}
\summaryhead

A three-paragraph news item. Marcus Einfeld, a retired Australian judge with a long list of honours, went to jail for two years for perjury and lies that began with a \$77 speeding ticket. The author presents the case as a real example of the theory of ``Entangled Truths, Contagious Lies'': that lies leave traces and can blow up in a rare, catastrophic way. The author adds that the theory sounds ``suspiciously virtuous'' and that selective reporting explains why we hear of this case.

\responsehead

The facts are right as far as they were known when the post appeared. The honours, the \$77 ticket for 6 mph over the limit, and the two years in jail all match the Wikipedia article the post links; only the claim that he was ``routinely brought back to judge important cases'' could not be confirmed.

The post's caveats are also right, and they limit what the case can show. A case that reaches us because it is remarkable cannot show how often lies blow up. It can show that they sometimes do, which ``Entangled Truths, Contagious Lies'' had already granted. The post claims ``a certain number of exemplars'' and no more.

As an exemplar, the case fits one part of the theory better than the other. It fits the tangled web: to defend one false name, Einfeld invented a second woman of the same name and produced a 20-page statement about her, which became the basis of one of the two charges that sent him to prison. It fits ``honest people not being good at lying'' less well. Evidence put to the Court of Appeal in July 2009, after the post, was that he had named absent people as the drivers of his car on earlier occasions, going back to 1999, and Wikipedia reports that the court found this proven. Those declarations seem to have come to light only after the 2006 case. On that record the case is also an example of lies that went uncaught for years, until the same kind of lie was told again.

In short: an accurately reported case, offered with fair caveats that limit it to showing that lies sometimes blow up; later evidence suggests that the same kind of lie had earlier gone undetected for years.
\end{afterword}
